\subsection{PARABOLIC SUBGROUPS}

DEFINITION 2. A subgroup of G is said to be parabolic if it contains a conjugate of B.

It is clear that every subgroup that contains a parabolic subgroup is parabolic.

PROPOSITION 4. Let P be a subgroup of G.

\begin{enumerate}
\def\labelenumi{\alph{enumi})}
\item
  P is parabolic if and only if there exists a subset X of S such that P is conjugate to \(G_{X}\) (cf.~no. 5 for the definition of \(G_{X}\) ).
\item
  Let \(\mathrm{X}, \mathrm{X}' \subset \mathrm{S}\) and \(g, g' \in \mathrm{G}\) be such that \(\mathrm{P} = g\mathrm{G}_{\mathrm{X}}g^{-1} = g'\mathrm{G}_{\mathrm{X}'}g'^{-1}\) . Then, \(\mathrm{X} = \mathrm{X}'\) and \(g'g^{-1} \in \mathrm{P}\) .
\end{enumerate}

Assertion a) follows from Th. 3, b).

Under the hypotheses of \(b\) ), we have

\[
g ^ {- 1} g ^ {\prime} \mathrm{B} g ^ {\prime - 1} g \subset g ^ {- 1} g ^ {\prime} \mathrm{G} _ {\mathrm{X} ^ {\prime}} g ^ {\prime - 1} g = \mathrm{G} _ {\mathrm{X}},
\]

and Prop. 3 shows that \(g^{-1}g' \in \mathrm{G}_X\) . Hence, \(\mathrm{G}_{\mathrm{X}'} = \mathrm{G}_X\) and \(\mathrm{X}' = \mathrm{X}\) by Th. 3, b). Finally,

\[
g ^ {\prime} g ^ {- 1} = g. g ^ {- 1} g ^ {\prime}. g ^ {- 1} \in g \mathrm{G} _ {\mathrm{X}} g ^ {- 1},
\]

which proves b).

If the parabolic subgroup P is conjugate to \(G_{X}\) , where \(X \subset S\) , then P is said to be of type X.

THEOREM 4. (i) Let \(\mathbf{P}_1\) and \(\mathbf{P}_2\) be two parabolic subgroups of \(\mathbf{G}\) whose intersection is parabolic and let \(g \in \mathbf{G}\) be such that \(g\mathrm{P}_1g^{-1} \subset \mathrm{P}_2\) . Then \(g \in \mathbf{P}_2\) and \(\mathbf{P}_1 \subset \mathbf{P}_2\) .

\begin{enumerate}
\def\labelenumi{(\roman{enumi})}
\setcounter{enumi}{1}
\item
  Two parabolic subgroups whose intersection is parabolic are not conjugate.
\item
  Let \(\mathbf{Q}_1\) and \(\mathbf{Q}_2\) be two parabolic subgroups of \(\mathbf{G}\) contained in a subgroup \(\mathbf{Q}\) of \(\mathbf{G}\) . Then any \(g \in \mathbf{G}\) such that \(g\mathbf{Q}_1g^{-1} = \mathbf{Q}_2\) belongs to \(\mathbf{Q}\) .
\item
  Every parabolic subgroup is its own normaliser \(^{6}\) .
\end{enumerate}

Assertion (i) follows from Props. 3 and 4, and implies (ii). Under the hypotheses of (iii), we have \(gQ_{1}g^{-1} \subset Q\) , which implies that \(g \in Q\) by (i). Finally, (iv) follows from (iii) by taking \(Q_{1} = Q_{2} = Q\) .

PROPOSITION 5. Let \(\mathbf{P}_1\) and \(\mathbf{P}_2\) be two parabolic subgroups of \(\mathbf{G}\) . Then \(\mathbf{P}_1 \cap \mathbf{P}_2\) contains a conjugate of \(\mathbf{T}\) .

By first transforming \(P_{1}\) and \(P_{2}\) by an inner automorphism of G, we may assume that \(B \subset P_{1}\) . Let \(g \in G\) be such that \(gBg^{-1} \subset P_{2}\) . By Th. 1, there exist \(n \in N\) and \(b, b' \in B\) such that \(g = bnb'\) . Since T is normal in N,

\[
\mathrm{P} _ {2} \supset g \mathrm{B} g ^ {- 1} = b n \mathrm{B} n ^ {- 1} b ^ {- 1} \supset b n \mathrm{T} n ^ {- 1} b ^ {- 1} = b \mathrm{T} b ^ {- 1}
\]

and

\[
\mathrm{P} _ {1} \supset \mathrm{B} \supset b \mathrm{T} b ^ {- 1},
\]

which proves the proposition.

